% Source: G10_T1_C1C2C3C4C5C6C7_exam.tex (questions and assessment rules).
% Full-width worked solutions; monetary units are local to each problem.
\makeatletter
\def\input@path{{./}{../}{../../}{../../../}{../../../../}{preamble/}{../preamble/}{../../preamble/}{../../../preamble/}{../../../../preamble/}{book2026-2027/preamble/}{book2026-2027/}}
\makeatother
\input{preamble_G10_T1_C1C2C3}
\input{preamble_exam_layout}
% Local solution-document spacing and readable labels in the inherited table.
\linespread{1.08}
\setlength{\parskip}{4pt}
\title{Grade 10 -- Term 1\\C1--C7 Catch-up Exam Solutions}
\author{}
\date{2026--2027}
\begin{document}
\maketitle
\vspace{-1.5cm}
\setcounter{tocdepth}{1}
\setcounter{secnumdepth}{2}
\phantomsection\label{toc}
\tableofcontents
\clearpage
\begingroup\small
\section*{Learning objective}
Characterise economic decisions under uncertainty, organise payoff tables, calculate profits from revenues and costs, and justify choices using Maximax, Maximin, Minimax Regret, EMV/Bayes, and probability sensitivity analysis.

\section*{Criteria assessment}
This exam assesses C1--C7. Each question and task shows its points.
\textbf{Passing requires a minimum of 5 cumulative points per criterion,
including points earned in previous learning activities}.
Show all procedures; answers without the required procedure receive no credit.
Partial credit may be awarded for a correct method with arithmetic errors.
Missing names or late submissions incur a 1-point deduction per criterion.
\section*{Evaluated criteria}
\begin{enumerate}[label=\textbf{C\arabic*}]
\item Characterises decision-making problems under uncertainty by identifying decision alternatives, states of nature, consequences, and the economic or financial context of the decision.
\item Organises decision problems using payoff tables, placing alternatives, states of nature, and payoffs in a clear structure for comparison.
\item Calculates profits by combining revenues with fixed and variable costs across different states of nature, determining the corresponding payoffs for each alternative, and organising the results in payoff tables for comparison.
\item Applies the Maximin and Maximax criteria to make optimistic or conservative decisions without probabilities, selecting the alternative with the highest possible payoff or the best worst-case, interpreting risks or limitations.
\item Constructs opportunity loss tables and applies the Minimax Regret criterion to select the alternative that minimises the maximum opportunity loss.
\item Calculates Expected Monetary Value using stated or estimated probabilities and applies the Bayes decision criterion to select the alternative with the highest expected payoff.
\item Uses subjective probabilities and sensitivity analysis to examine how changes in probability assumptions affect the preferred decision.
\end{enumerate}

\begin{center}
	\small\setlength{\tabcolsep}{5pt}
	\begin{tabular}{|c|c|c|c|c||c|c|c|c|}
		\hline
		\textbf{\textnumero}&\textbf{C1}&\textbf{C2}&\textbf{C3}&\textbf{C4}&
		\textbf{\textnumero}&\textbf{C5}&\textbf{C6}&\textbf{C7}\\\hline
		1&3&6&--&--&6&2&--&--\\\hline
		2&3&--&6&--&7&2&--&--\\\hline
		3&--&--&--&2&8&2&--&--\\\hline
		4&--&--&--&2&9&--&2&--\\\hline
		5&--&--&--&2&10&--&2&--\\\hline
		&&&&&11&--&2&--\\\hline
		&&&&&12&--&--&3\\\hline
		&&&&&13&--&--&3\\\hline
		\textbf{Tot.}&\textbf{6}&\textbf{6}&\textbf{6}&\textbf{6}&
		\textbf{Tot.}&\textbf{6}&\textbf{6}&\textbf{6}\\\hline
	\end{tabular}
\end{center}

Each criterion has 6 available points; the exam has 42 available points overall.
\textbf{Notation.} $A_i$ and $S_j$ always follow the original alternative and state order within each problem. $P(A_i,S_j)$ denotes a monetary payoff; $P(S_j)$ denotes a probability. Navigation arrows return to the contents, the problem, or the summary.
\endgroup

\clearpage
\sectionScored{Craft Workshop}{C1-\sml{3} C2-\sml{6}}{sec:craft-workshop}
A craft studio must choose a pottery workshop or a weaving workshop before tourist demand is known. Demand will be strong or weak. Pottery provides a net contribution of 12 USD per participant and a fixed sponsorship contribution of 80 USD per workshop. Weaving provides 15 USD per participant and a fixed sponsorship contribution of 50 USD per workshop. Strong demand adds a further net contribution of 6 USD per participant to either workshop; weak demand adds 2 USD per participant. Pottery attracts 30 participants with strong demand or 18 with weak demand; weaving attracts 24 or 16, respectively. All contributions are additive, and no other adjustments are needed.
\QuestionsTasks[sec:craft-workshop]{
{1}/{sol:craft-workshop-c1}/{Characterize the key elements of the decision-making situation \sscore{C1-3}},
{2}/{sol:craft-workshop-c2}/{and construct the payoff table. \sscore{C2-6}}
}
\subsection*{C1 - Characterization \score{3}}
\phantomsection\label{sol:craft-workshop-c1}
\DecisionElementsTable{Pottery workshop; Weaving workshop}{Strong tourist demand; Weak tourist demand}{Tourist demand revealed after the workshop is chosen; the studio cannot control it.}{Operating the selected activity under the realised conditions leaves funds available or a shortfall to cover. Monetary measure: Total net contribution in USD for each workshop--demand pair.}
The economic context is a craft studio choosing its workshop before tourist demand is known.
The decision controls the workshop, while demand controls attendance and the extra contribution rate.
\begin{CalculationLandscape}{11in}{8.5in}
\subsection*{C2 - Additive contributions and payoff table \score{6}}
\phantomsection\label{sol:craft-workshop-c2}

\SymbolsTable{
$A_i$ & $A_1$: Pottery workshop; $A_2$:  Weaving workshop & Alternative labels, in row order. \\
$S_j$ & $S_1$: Strong tourist demand; $S_2$:  Weak tourist demand & State labels, in column order; revealed after the choice. \\
$q_{ij}$ & $q_{11}=30$ participants; $q_{12}=18$ participants; $q_{21}=24$ participants; $q_{22}=16$ participants & Attendance for alternative $A_i$ under state $S_j$. \\
$v_i$ & 12; 15 USD/participant & Alternative-based net contribution, in row order. \\
$d_j$ & 6; 2 USD/participant & State-based extra net contribution, in column order. \\
$f_i$ & 80; 50 USD & Fixed sponsorship per chosen activity, in row order. \\
}

Let $q_{ij}$ be participants, $v_i$ the workshop contribution per participant,
$d_j$ the demand contribution per participant, and $f_i$ the fixed sponsorship.
All monetary rates are in USD per participant; $f_i$ and payoffs are in USD.
\[
P_{ay}(A_i,S_j)=q_{ij}(v_i+d_j)+f_i.
\]
The per-participant contributions are added before multiplying by participants.
The fixed sponsorship is then added once per workshop; it is not multiplied by participants.
\CalculationPair{Additive payoff (USD)}{
\begin{tabular}{@{}lcc@{}}
\toprule Alternative & $S_1$ & $S_2$ \\ \midrule
$A_1$ & \SymbolicCell{q_{11}\times \\ [v_1+d_1] \\ {}+f_1}{P_{ay}(A_1,S_1)} & \SymbolicCell{q_{12}\times \\ [v_1+d_2] \\ {}+f_1}{P_{ay}(A_1,S_2)} \\[5pt]
$A_2$ & \SymbolicCell{q_{21}\times \\ [v_2+d_1] \\ {}+f_2}{P_{ay}(A_2,S_1)} & \SymbolicCell{q_{22}\times \\ [v_2+d_2] \\ {}+f_2}{P_{ay}(A_2,S_2)} \\[5pt]
\bottomrule
\end{tabular}
}{
\begin{tabular}{@{}lcc@{}}
\toprule Alternative & $S_1$ & $S_2$ \\ \midrule
$A_1$ & \CalcResult{30\times \\ [12+6] \\ {}+80}{620} & \CalcResult{18\times \\ [12+2] \\ {}+80}{332} \\[5pt]
$A_2$ & \CalcResult{24\times \\ [15+6] \\ {}+50}{554} & \CalcResult{16\times \\ [15+2] \\ {}+50}{322} \\[5pt]
\bottomrule
\end{tabular}
}

\Needspace{10\baselineskip}
\textbf{Completed payoff table (USD).}
\begin{center}
\begin{tabular}{@{}lcc@{}}\toprule
Alternative & Strong ($S_1$) & Weak ($S_2$)\\\midrule
Pottery workshop ($A_1$) & 620 & 332\\
Weaving workshop ($A_2$) & 554 & 322\\
\bottomrule\end{tabular}
\end{center}
\end{CalculationLandscape}

\textbf{Conclusion:} the pottery workshop gives a higher contribution in both states: USD 620 under strong demand and USD 332 under weak demand.
\secInicio[sec:craft-workshop][sec:summarized-solutions]

\clearpage
\sectionScored{Guided Cycling Tours}{C1-\sml{3} C3-\sml{6}}{sec:guided-cycling-tours}
\begingroup\setlength{\parskip}{2pt}\linespread{1.02}\selectfont
A tour operator must choose a town route or a countryside route before seasonal demand is known. Demand will be strong or weak. The town route sells 90 places with strong demand or 60 with weak demand; the countryside route sells 70 or 40, respectively. Base revenue per place is 32 USD for the town route and 45 USD for the countryside route. Strong demand adds a surcharge of 8 USD per place to either route; weak demand adds 3 USD per place. Variable cost per place is 18 USD for the town route and 26 USD for the countryside route. Total fixed seasonal costs, for either route, are 500 USD with strong demand and 650 USD with weak demand. These fixed costs are paid once per season, and there are no other revenues or costs.
\QuestionsTasks[sec:guided-cycling-tours]{
{1}/{sol:guided-cycling-tours-c1}/{Characterize the key elements of the decision-making situation \sscore{C1-3}},
{3}/{sol:guided-cycling-tours-c3}/{and construct the payoff table. \sscore{C3-6}}
}
\subsection*{C1 - Characterization \score{3}}
\phantomsection\label{sol:guided-cycling-tours-c1}
\DecisionElementsTable{Town route; Countryside route}{Strong seasonal demand; Weak seasonal demand}{Seasonal demand revealed after the route is chosen; the operator cannot control it.}{Operating the selected activity under the realised conditions leaves funds available or a shortfall to cover. Monetary measure: Seasonal profit in USD for each route--demand pair.}
The economic context is a tour operator choosing a route before knowing sales,
the demand surcharge, and the seasonal fixed cost.
\begin{CalculationLandscape}{14in}{11in}
\subsection*{C3 - Profit calculations and payoff table \score{6}}
\phantomsection\label{sol:guided-cycling-tours-c3}

\SymbolsTable{
$A_i$ & $A_1$: Town route; $A_2$:  Countryside route & Alternative labels, in row order. \\
$S_j$ & $S_1$: Strong seasonal demand; $S_2$:  Weak seasonal demand & State labels, in column order; revealed after the choice. \\
$q_{ij}$ & $q_{11}=90$ places; $q_{12}=60$ places; $q_{21}=70$ places; $q_{22}=40$ places & Sales quantity for alternative $A_i$ under state $S_j$. \\
$b_i$ & 32; 45 USD/place & Base revenue per sale, in alternative order. \\
$s_j$ & 8; 3 USD/place & Revenue surcharge per sale, in state order. \\
$c_i$ & 18; 26 USD/place & Variable cost per sale, in alternative order. \\
$F_j$ & 500; 650 USD & Fixed seasonal cost, paid once; in state order. \\
}

\[C(A_i,S_j)=VC_{ij}+F_j,\qquad P(A_i,S_j)=R(A_i,S_j)-C(A_i,S_j).\]
Profit equals total revenue minus total variable costs minus fixed seasonal costs:
\begin{align*}
R_{ij}&=q_{ij}(b_i+s_j), & VC_{ij}&=q_{ij}c_i,\\
P(A_i,S_j)&=R_{ij}-VC_{ij}-F_j
=q_{ij}(b_i+s_j)-q_{ij}c_i-F_j.
\end{align*}
Here $q_{ij}$ is places sold; $b_i$, $s_j$, and $c_i$ are base revenue, surcharge,
and variable cost in USD per place; $F_j$ is the fixed seasonal cost in USD.
The surcharge is part of revenue for every place sold. The fixed cost is paid
once in each alternative--state calculation, not once per place.
\CalculationPair{Revenue (USD)}{
\begin{tabular}{@{}lcc@{}}
\toprule Alternative & $S_1$ & $S_2$ \\ \midrule
$A_1$ & \SymbolicCell{q(A_1,S_1)\times \\ [b_1+s_1]}{R(A_1,S_1)} & \SymbolicCell{q(A_1,S_2)\times \\ [b_1+s_2]}{R(A_1,S_2)} \\[5pt]
$A_2$ & \SymbolicCell{q(A_2,S_1)\times \\ [b_2+s_1]}{R(A_2,S_1)} & \SymbolicCell{q(A_2,S_2)\times \\ [b_2+s_2]}{R(A_2,S_2)} \\[5pt]
\bottomrule
\end{tabular}
}{
\begin{tabular}{@{}lcc@{}}
\toprule Alternative & $S_1$ & $S_2$ \\ \midrule
$A_1$ & \CalcResult{90\times \\ [32+8]}{3600} & \CalcResult{60\times \\ [32+3]}{2100} \\[5pt]
$A_2$ & \CalcResult{70\times \\ [45+8]}{3710} & \CalcResult{40\times \\ [45+3]}{1920} \\[5pt]
\bottomrule
\end{tabular}
}
\CalculationPair{Cost (USD)}{
\begin{tabular}{@{}lcc@{}}
\toprule Alternative & $S_1$ & $S_2$ \\ \midrule
$A_1$ & \SymbolicCell{q(A_1,S_1)\times c_1 \\ {}+F_1}{C(A_1,S_1)} & \SymbolicCell{q(A_1,S_2)\times c_1 \\ {}+F_2}{C(A_1,S_2)} \\[5pt]
$A_2$ & \SymbolicCell{q(A_2,S_1)\times c_2 \\ {}+F_1}{C(A_2,S_1)} & \SymbolicCell{q(A_2,S_2)\times c_2 \\ {}+F_2}{C(A_2,S_2)} \\[5pt]
\bottomrule
\end{tabular}
}{
\begin{tabular}{@{}lcc@{}}
\toprule Alternative & $S_1$ & $S_2$ \\ \midrule
$A_1$ & \CalcResult{90\times 18 \\ {}+500}{2120} & \CalcResult{60\times 18 \\ {}+650}{1730} \\[5pt]
$A_2$ & \CalcResult{70\times 26 \\ {}+500}{2320} & \CalcResult{40\times 26 \\ {}+650}{1690} \\[5pt]
\bottomrule
\end{tabular}
}
\CalculationPair{Profit (USD)}{
\begin{tabular}{@{}lcc@{}}
\toprule Alternative & $S_1$ & $S_2$ \\ \midrule
$A_1$ & \SymbolicCell{R(A_1,S_1)-C(A_1,S_1)}{P(A_1,S_1)} & \SymbolicCell{R(A_1,S_2)-C(A_1,S_2)}{P(A_1,S_2)} \\[5pt]
$A_2$ & \SymbolicCell{R(A_2,S_1)-C(A_2,S_1)}{P(A_2,S_1)} & \SymbolicCell{R(A_2,S_2)-C(A_2,S_2)}{P(A_2,S_2)} \\[5pt]
\bottomrule
\end{tabular}
}{
\begin{tabular}{@{}lcc@{}}
\toprule Alternative & $S_1$ & $S_2$ \\ \midrule
$A_1$ & \CalcResult{3600-2120}{1480} & \CalcResult{2100-1730}{370} \\[5pt]
$A_2$ & \CalcResult{3710-2320}{1390} & \CalcResult{1920-1690}{230} \\[5pt]
\bottomrule
\end{tabular}
}

\Needspace{10\baselineskip}
\textbf{Completed payoff table (profit in USD).}
\begin{center}
\begin{tabular}{@{}lcc@{}}\toprule
Alternative & Strong ($S_1$) & Weak ($S_2$)\\\midrule
Town route ($A_1$) & 1480 & 370\\
Countryside route ($A_2$) & 1390 & 230\\
\bottomrule\end{tabular}
\end{center}
\end{CalculationLandscape}

\textbf{Conclusion:} the town route earns more in both states: USD 1,480 under strong demand and USD 370 under weak demand.
\secInicio[sec:guided-cycling-tours][sec:summarized-solutions]
\endgroup

\clearpage
\sectionScored{Sportswear Stock}{C4-\sml{2}}{sec:sportswear-stock}
A retailer must choose a seasonal stock plan before demand is known. Demand will be high, medium, or low. Profits are in thousand USD; no probabilities are available.
\begin{center}\small\setlength{\tabcolsep}{4pt}
\begin{tabular}{@{}lrrr@{}}\toprule
Alternative&High&Medium&Low\\\midrule
Bulk order&120&48&-30\\
Regular order&84&57&16\\
Small order&55&43&29\\
Pre-orders only&24&24&24\\\bottomrule
\end{tabular}\end{center}
\QuestionsTasks[sec:sportswear-stock]{
{4}/{sol:sportswear-stock-c4}/{What would the optimistic decision be? What would the conservative decision be? Justify your answer. \score{2}}
}
\subsection*{C4 - Maximax and Maximin \score{2}}
\phantomsection\label{sol:sportswear-stock-c4}
Alternatives: $A_1$: Bulk order; $A_2$: Regular order; $A_3$: Small order; $A_4$: Pre-orders only.\par
States: $S_1$: High; $S_2$: Medium; $S_3$: Low.
All profits below are in thousand USD. For every row, find the largest and smallest payoff:
\begin{center}
\begin{tabular}{@{}lcc@{}}\toprule
Alternative & Row maximum $M_i$ & Row minimum $m_i$\\\midrule
Bulk order ($A_1$) & $\max\{120,48,-30\}=120$ & $\min\{120,48,-30\}=-30$\\
Regular order ($A_2$) & $\max\{84,57,16\}=84$ & $\min\{84,57,16\}=16$\\
Small order ($A_3$) & $\max\{55,43,29\}=55$ & $\min\{55,43,29\}=29$\\
Pre-orders only ($A_4$) & $\max\{24,24,24\}=24$ & $\min\{24,24,24\}=24$\\
\bottomrule\end{tabular}
\end{center}
\[\max_i M_i=\max\{120,84,55,24\}=120,\qquad \max_i m_i=\max\{-30,16,29,24\}=29.\]
\textbf{Maximax:} choose \textbf{Bulk order} ($A_1$), with best-case profit 120 thousand USD.\par
\textbf{Maximin:} choose \textbf{Small order} ($A_3$), with worst-case profit 29 thousand USD.
Maximax is optimistic: it focuses on the best possible result and ignores the downside. Maximin is conservative: it protects the best worst-case result and may forgo a large gain. Without probabilities, neither method estimates how likely a payoff is or gives an expected profit.
\secInicio[sec:sportswear-stock][sec:summarized-solutions]

\clearpage
\sectionScored{Concert Venue}{C4-\sml{2}}{sec:concert-venue}
An organiser must book one venue before ticket demand is known. Demand will be high, medium, or low. Profits are in thousand USD; no probabilities are available.
\begin{center}\small\setlength{\tabcolsep}{4pt}
\begin{tabular}{@{}lrrr@{}}\toprule
Alternative&High&Medium&Low\\\midrule
Stadium&150&42&-45\\
Theatre&95&62&12\\
Community hall&60&46&26\\
Online event&30&30&30\\\bottomrule
\end{tabular}\end{center}
\QuestionsTasks[sec:concert-venue]{
{4}/{sol:concert-venue-c4}/{What would the optimistic decision be? What would the conservative decision be? Justify your answer. \score{2}}
}
\subsection*{C4 - Maximax and Maximin \score{2}}
\phantomsection\label{sol:concert-venue-c4}
Alternatives: $A_1$: Stadium; $A_2$: Theatre; $A_3$: Community hall; $A_4$: Online event.\par
States: $S_1$: High; $S_2$: Medium; $S_3$: Low.
All profits below are in thousand USD. For every row, find the largest and smallest payoff:
\begin{center}
\begin{tabular}{@{}lcc@{}}\toprule
Alternative & Row maximum $M_i$ & Row minimum $m_i$\\\midrule
Stadium ($A_1$) & $\max\{150,42,-45\}=150$ & $\min\{150,42,-45\}=-45$\\
Theatre ($A_2$) & $\max\{95,62,12\}=95$ & $\min\{95,62,12\}=12$\\
Community hall ($A_3$) & $\max\{60,46,26\}=60$ & $\min\{60,46,26\}=26$\\
Online event ($A_4$) & $\max\{30,30,30\}=30$ & $\min\{30,30,30\}=30$\\
\bottomrule\end{tabular}
\end{center}
\[\max_i M_i=\max\{150,95,60,30\}=150,\qquad \max_i m_i=\max\{-45,12,26,30\}=30.\]
\textbf{Maximax:} choose \textbf{Stadium} ($A_1$), with best-case profit 150 thousand USD.\par
\textbf{Maximin:} choose \textbf{Online event} ($A_4$), with worst-case profit 30 thousand USD.
Maximax is optimistic: it focuses on the best possible result and ignores the downside. Maximin is conservative: it protects the best worst-case result and may forgo a large gain. Without probabilities, neither method estimates how likely a payoff is or gives an expected profit.
\secInicio[sec:concert-venue][sec:summarized-solutions]

\clearpage
\sectionScored{Cold Storage}{C4-\sml{2}}{sec:cold-storage}
A food distributor must choose a storage plan before next season's market conditions are known. The four possible states are S1: very strong sales, S2: strong sales, S3: steady sales, and S4: weak sales. Profits are in thousand USD; no probabilities are available.
\begin{center}\small\setlength{\tabcolsep}{4pt}
\begin{tabular}{@{}lrrrr@{}}\toprule
Alternative&S1&S2&S3&S4\\\midrule
Build large&140&72&14&-40\\
Build medium&100&76&38&8\\
Rent small&68&56&42&22\\
Shared storage&48&44&36&28\\\bottomrule
\end{tabular}\end{center}
\QuestionsTasks[sec:cold-storage]{
{4}/{sol:cold-storage-c4}/{What would the optimistic decision be? What would the conservative decision be? Justify your answer. \score{2}}
}
\subsection*{C4 - Maximax and Maximin \score{2}}
\phantomsection\label{sol:cold-storage-c4}
Alternatives: $A_1$: Build large; $A_2$: Build medium; $A_3$: Rent small; $A_4$: Shared storage.\par
States: $S_1$: Very strong sales; $S_2$: Strong sales; $S_3$: Steady sales; $S_4$: Weak sales.
All profits below are in thousand USD. For every row, find the largest and smallest payoff:
\begin{center}
\begin{tabular}{@{}lcc@{}}\toprule
Alternative & Row maximum $M_i$ & Row minimum $m_i$\\\midrule
Build large ($A_1$) & $\max\{140,72,14,-40\}=140$ & $\min\{140,72,14,-40\}=-40$\\
Build medium ($A_2$) & $\max\{100,76,38,8\}=100$ & $\min\{100,76,38,8\}=8$\\
Rent small ($A_3$) & $\max\{68,56,42,22\}=68$ & $\min\{68,56,42,22\}=22$\\
Shared storage ($A_4$) & $\max\{48,44,36,28\}=48$ & $\min\{48,44,36,28\}=28$\\
\bottomrule\end{tabular}
\end{center}
\[\max_i M_i=\max\{140,100,68,48\}=140,\qquad \max_i m_i=\max\{-40,8,22,28\}=28.\]
\textbf{Maximax:} choose \textbf{Build large} ($A_1$), with best-case profit 140 thousand USD.\par
\textbf{Maximin:} choose \textbf{Shared storage} ($A_4$), with worst-case profit 28 thousand USD.
Maximax is optimistic: it focuses on the best possible result and ignores the downside. Maximin is conservative: it protects the best worst-case result and may forgo a large gain. Without probabilities, neither method estimates how likely a payoff is or gives an expected profit.
\secInicio[sec:cold-storage][sec:summarized-solutions]

\clearpage
\sectionScored{Flower Stall}{C5-\sml{2}}{sec:flower-stall}
A florist must choose a purchasing plan before holiday demand is known. The three possible demand states are high, medium, and low. Profits are in thousand USD; no probabilities are available. All payoffs are to be maximised.
\begin{center}\small\setlength{\tabcolsep}{4pt}
\begin{tabular}{@{}lrrr@{}}\toprule
Alternative&High&Medium&Low\\\midrule
Large purchase&78&32&-12\\
Medium purchase&60&48&18\\
Small purchase&40&38&30\\\bottomrule
\end{tabular}\end{center}
\QuestionsTasks[sec:flower-stall]{
{5}/{sol:flower-stall-c5}/{What decision would be selected using the Minimax Regret criterion? Justify your answer. \score{2}}
}
\subsection*{C5 - Minimax Regret \score{2}}
\phantomsection\label{sol:flower-stall-c5}
Alternatives: $A_1$: Large purchase; $A_2$: Medium purchase; $A_3$: Small purchase.\par
States: $S_1$: High; $S_2$: Medium; $S_3$: Low.
\[\boldsymbol{b}=\bigl(78,48,30\bigr),\qquad b_j=\max_k P(A_k,S_j).\]
Each column is compared with its highest profit. Regret is the profit missed by not choosing that state's best alternative:
\[r(A_i,S_j)=\max_k P(A_k,S_j)-P(A_i,S_j).\]
All payoffs and regrets in this calculation are in thousand USD.
\begin{center}
\begin{tabular}{@{}lcccc@{}}\toprule
Alternative & $S_1$ & $S_2$ & $S_3$ & Maximum regret\\\midrule
$A_1$ & $\begin{gathered}78-78\\=0\end{gathered}$ & $\begin{gathered}48-32\\=16\end{gathered}$ & $\begin{gathered}30-(-12)\\=42\end{gathered}$ & 42\\
$A_2$ & $\begin{gathered}78-60\\=18\end{gathered}$ & $\begin{gathered}48-48\\=0\end{gathered}$ & $\begin{gathered}30-18\\=12\end{gathered}$ & 18\\
$A_3$ & $\begin{gathered}78-40\\=38\end{gathered}$ & $\begin{gathered}48-38\\=10\end{gathered}$ & $\begin{gathered}30-30\\=0\end{gathered}$ & 38\\
\bottomrule\end{tabular}
\end{center}
\[\min_i\max_j r(A_i,S_j)=\min\{42,18,38\}=18.\]
Choose \textbf{Medium purchase} ($A_2$): its maximum regret is \textbf{18 thousand USD}. This is the smallest worst possible opportunity loss; it does not mean that this loss occurs in every state or that this alternative always earns the most profit.
\secInicio[sec:flower-stall][sec:summarized-solutions]

\clearpage
\sectionScored{Parcel Depot}{C5-\sml{2}}{sec:parcel-depot}
A parcel company must choose a depot plan before shipment demand is known. Demand will be high, normal, or low. Profits are in thousand USD; no probabilities are available. All payoffs are to be maximised.
\begin{center}\small\setlength{\tabcolsep}{4pt}
\begin{tabular}{@{}lrrr@{}}\toprule
Alternative&High&Normal&Low\\\midrule
New depot&112&40&-28\\
Expand depot&86&64&14\\
Rent space&62&54&34\\
Current depot&36&36&36\\\bottomrule
\end{tabular}\end{center}
\QuestionsTasks[sec:parcel-depot]{
{5}/{sol:parcel-depot-c5}/{What decision would be selected using the Minimax Regret criterion? Justify your answer. \score{2}}
}
\subsection*{C5 - Minimax Regret \score{2}}
\phantomsection\label{sol:parcel-depot-c5}
Alternatives: $A_1$: New depot; $A_2$: Expand depot; $A_3$: Rent space; $A_4$: Current depot.\par
States: $S_1$: High; $S_2$: Normal; $S_3$: Low.
\[\boldsymbol{b}=\bigl(112,64,36\bigr),\qquad b_j=\max_k P(A_k,S_j).\]
Each column is compared with its highest profit. Regret is the profit missed by not choosing that state's best alternative:
\[r(A_i,S_j)=\max_k P(A_k,S_j)-P(A_i,S_j).\]
All payoffs and regrets in this calculation are in thousand USD.
\begin{center}
\begin{tabular}{@{}lcccc@{}}\toprule
Alternative & $S_1$ & $S_2$ & $S_3$ & Maximum regret\\\midrule
$A_1$ & $\begin{gathered}112-112\\=0\end{gathered}$ & $\begin{gathered}64-40\\=24\end{gathered}$ & $\begin{gathered}36-(-28)\\=64\end{gathered}$ & 64\\
$A_2$ & $\begin{gathered}112-86\\=26\end{gathered}$ & $\begin{gathered}64-64\\=0\end{gathered}$ & $\begin{gathered}36-14\\=22\end{gathered}$ & 26\\
$A_3$ & $\begin{gathered}112-62\\=50\end{gathered}$ & $\begin{gathered}64-54\\=10\end{gathered}$ & $\begin{gathered}36-34\\=2\end{gathered}$ & 50\\
$A_4$ & $\begin{gathered}112-36\\=76\end{gathered}$ & $\begin{gathered}64-36\\=28\end{gathered}$ & $\begin{gathered}36-36\\=0\end{gathered}$ & 76\\
\bottomrule\end{tabular}
\end{center}
\[\min_i\max_j r(A_i,S_j)=\min\{64,26,50,76\}=26.\]
Choose \textbf{Expand depot} ($A_2$): its maximum regret is \textbf{26 thousand USD}. This is the smallest worst possible opportunity loss; it does not mean that this loss occurs in every state or that this alternative always earns the most profit.
\secInicio[sec:parcel-depot][sec:summarized-solutions]

\clearpage
\sectionScored{Bakery Output}{C5-\sml{2}}{sec:bakery-output}
A bakery must choose its daily output before demand is known. The four possible demand states are S1: very high, S2: high, S3: moderate, and S4: low. Profits are in USD; no probabilities are available. All payoffs are to be maximised.
\begin{center}\small\setlength{\tabcolsep}{4pt}
\begin{tabular}{@{}lrrrr@{}}\toprule
Alternative&S1&S2&S3&S4\\\midrule
High output&960&560&180&-240\\
Flexible output&740&660&460&120\\
Low output&480&500&440&360\\\bottomrule
\end{tabular}\end{center}
\QuestionsTasks[sec:bakery-output]{
{5}/{sol:bakery-output-c5}/{What decision would be selected using the Minimax Regret criterion? Justify your answer. \score{2}}
}
\subsection*{C5 - Minimax Regret \score{2}}
\phantomsection\label{sol:bakery-output-c5}
Alternatives: $A_1$: High output; $A_2$: Flexible output; $A_3$: Low output.\par
States: $S_1$: Very high; $S_2$: High; $S_3$: Moderate; $S_4$: Low.
\[\boldsymbol{b}=\bigl(960,660,460,360\bigr),\qquad b_j=\max_k P(A_k,S_j).\]
Each column is compared with its highest profit. Regret is the profit missed by not choosing that state's best alternative:
\[r(A_i,S_j)=\max_k P(A_k,S_j)-P(A_i,S_j).\]
All payoffs and regrets in this calculation are in USD.
\begin{center}
\begin{tabular}{@{}lccccc@{}}\toprule
Alternative & $S_1$ & $S_2$ & $S_3$ & $S_4$ & Maximum regret\\\midrule
$A_1$ & $\begin{gathered}960-960\\=0\end{gathered}$ & $\begin{gathered}660-560\\=100\end{gathered}$ & $\begin{gathered}460-180\\=280\end{gathered}$ & $\begin{gathered}360-(-240)\\=600\end{gathered}$ & 600\\
$A_2$ & $\begin{gathered}960-740\\=220\end{gathered}$ & $\begin{gathered}660-660\\=0\end{gathered}$ & $\begin{gathered}460-460\\=0\end{gathered}$ & $\begin{gathered}360-120\\=240\end{gathered}$ & 240\\
$A_3$ & $\begin{gathered}960-480\\=480\end{gathered}$ & $\begin{gathered}660-500\\=160\end{gathered}$ & $\begin{gathered}460-440\\=20\end{gathered}$ & $\begin{gathered}360-360\\=0\end{gathered}$ & 480\\
\bottomrule\end{tabular}
\end{center}
\[\min_i\max_j r(A_i,S_j)=\min\{600,240,480\}=240.\]
Choose \textbf{Flexible output} ($A_2$): its maximum regret is \textbf{240 USD}. This is the smallest worst possible opportunity loss; it does not mean that this loss occurs in every state or that this alternative always earns the most profit.
\secInicio[sec:bakery-output][sec:summarized-solutions]

\clearpage
\sectionScored{Fruit Packing}{C6-\sml{2}}{sec:fruit-packing}
A cooperative must choose a packing plan before export demand is known. Demand will be strong, steady, or weak, with the probabilities shown. Profits are in thousand USD.
\begin{center}\small\setlength{\tabcolsep}{4pt}
\begin{tabular}{@{}lrrr@{}}\toprule
Alternative&Strong&Steady&Weak\\
$P(S_j)$&0.40&0.35&0.25\\\midrule
Premium packs&90&44&-12\\
Mixed packs&68&58&24\\
Basic packs&46&48&38\\\bottomrule
\end{tabular}\end{center}
\QuestionsTasks[sec:fruit-packing]{
{6}/{sol:fruit-packing-c6}/{What decision would be selected using the EMV or Bayes criterion? Justify your answer. \score{2}}
}
\subsection*{C6 - EMV / Bayes \score{2}}
\phantomsection\label{sol:fruit-packing-c6}
Alternatives: $A_1$: Premium packs; $A_2$: Mixed packs; $A_3$: Basic packs.\par
States: $S_1$: Strong; $S_2$: Steady; $S_3$: Weak.
\[\boldsymbol{q}=(0.40,0.35,0.25),\qquad 0.40+0.35+0.25=1.\]
All probabilities lie in $[0,1]$, so they are nonnegative and valid.
\[EMV(A_i)=\sum_{j=1}^{n}P(S_j)P(A_i,S_j).\]
Multiply each state payoff by its probability, including any negative payoff, and add the weighted amounts. All EMVs are in thousand USD.
\begin{align*}
EMV(A_1)&=0.40(90)+0.35(44)+0.25(-12)\\&=36+15.40-3=48.40,\\
EMV(A_2)&=0.40(68)+0.35(58)+0.25(24)\\&=27.20+20.30+6=53.50,\\
EMV(A_3)&=0.40(46)+0.35(48)+0.25(38)\\&=18.40+16.80+9.50=44.70.
\end{align*}
\[\max\{48.40,53.50,44.70\}=53.50.\]
Choose \textbf{Mixed packs} ($A_2$), with expected profit \textbf{53.50 thousand USD} (USD 53,500). EMV is a probability-weighted expected payoff, not a guaranteed outcome. The realised profit will be the payoff for the state that actually occurs.
\secInicio[sec:fruit-packing][sec:summarized-solutions]

\clearpage
\sectionScored{Art Market Space}{C6-\sml{2}}{sec:art-market-space}
An art-market organiser must choose a space before visitor attendance is known. Attendance will be high, normal, or low, with the probabilities shown. Profits are in thousand USD.
\begin{center}\small\setlength{\tabcolsep}{4pt}
\begin{tabular}{@{}lrrr@{}}\toprule
Alternative&High&Normal&Low\\
$P(S_j)$&0.30&0.40&0.30\\\midrule
Large tent&104&38&-24\\
Medium tent&78&56&18\\
Small tent&56&50&36\\
Existing rooms&28&28&28\\\bottomrule
\end{tabular}\end{center}
\QuestionsTasks[sec:art-market-space]{
{6}/{sol:art-market-space-c6}/{What decision would be selected using the EMV or Bayes criterion? Justify your answer. \score{2}}
}
\subsection*{C6 - EMV / Bayes \score{2}}
\phantomsection\label{sol:art-market-space-c6}
Alternatives: $A_1$: Large tent; $A_2$: Medium tent; $A_3$: Small tent; $A_4$: Existing rooms.\par
States: $S_1$: High; $S_2$: Normal; $S_3$: Low.
\[\boldsymbol{q}=(0.30,0.40,0.30),\qquad 0.30+0.40+0.30=1.\]
All probabilities lie in $[0,1]$, so they are nonnegative and valid.
\[EMV(A_i)=\sum_{j=1}^{n}P(S_j)P(A_i,S_j).\]
Multiply each state payoff by its probability, including any negative payoff, and add the weighted amounts. All EMVs are in thousand USD.
\begin{align*}
EMV(A_1)&=0.30(104)+0.40(38)+0.30(-24)\\&=31.20+15.20-7.20=39.20,\\
EMV(A_2)&=0.30(78)+0.40(56)+0.30(18)\\&=23.40+22.40+5.40=51.20,\\
EMV(A_3)&=0.30(56)+0.40(50)+0.30(36)\\&=16.80+20+10.80=47.60,\\
EMV(A_4)&=0.30(28)+0.40(28)+0.30(28)\\&=8.40+11.20+8.40=28.
\end{align*}
\[\max\{39.20,51.20,47.60,28\}=51.20.\]
Choose \textbf{Medium tent} ($A_2$), with expected profit \textbf{51.20 thousand USD} (USD 51,200). EMV is a probability-weighted expected payoff, not a guaranteed outcome. The realised profit will be the payoff for the state that actually occurs.
\secInicio[sec:art-market-space][sec:summarized-solutions]

\clearpage
\sectionScored{Repair Shop Staffing}{C6-\sml{2}}{sec:repair-shop-staffing}
A repair shop must choose a staffing plan before service demand is known. The four possible demand states are S1: surge, S2: strong, S3: steady, and S4: slow, with the probabilities shown. Profits are in thousand USD.
\begin{center}\small\setlength{\tabcolsep}{4pt}
\begin{tabular}{@{}lrrrr@{}}\toprule
Alternative&S1&S2&S3&S4\\
$P(S_j)$&0.20&0.30&0.30&0.20\\\midrule
Hire full team&120&72&24&-36\\
Hire part-time&88&76&52&12\\
Current team&58&60&54&44\\\bottomrule
\end{tabular}\end{center}
\QuestionsTasks[sec:repair-shop-staffing]{
{6}/{sol:repair-shop-staffing-c6}/{What decision would be selected using the EMV or Bayes criterion? Justify your answer. \score{2}}
}
\subsection*{C6 - EMV / Bayes \score{2}}
\phantomsection\label{sol:repair-shop-staffing-c6}
Alternatives: $A_1$: Hire full team; $A_2$: Hire part-time; $A_3$: Current team.\par
States: $S_1$: Surge; $S_2$: Strong; $S_3$: Steady; $S_4$: Slow.
\[\boldsymbol{q}=(0.20,0.30,0.30,0.20),\qquad 0.20+0.30+0.30+0.20=1.\]
All probabilities lie in $[0,1]$, so they are nonnegative and valid.
\[EMV(A_i)=\sum_{j=1}^{n}P(S_j)P(A_i,S_j).\]
Multiply each state payoff by its probability, including any negative payoff, and add the weighted amounts. All EMVs are in thousand USD.
\begin{align*}
EMV(A_1)&=0.20(120)+0.30(72)+0.30(24)+0.20(-36)\\&=24+21.60+7.20-7.20=45.60,\\
EMV(A_2)&=0.20(88)+0.30(76)+0.30(52)+0.20(12)\\&=17.60+22.80+15.60+2.40=58.40,\\
EMV(A_3)&=0.20(58)+0.30(60)+0.30(54)+0.20(44)\\&=11.60+18+16.20+8.80=54.60.
\end{align*}
\[\max\{45.60,58.40,54.60\}=58.40.\]
Choose \textbf{Hire part-time} ($A_2$), with expected profit \textbf{58.40 thousand USD} (USD 58,400). EMV is a probability-weighted expected payoff, not a guaranteed outcome. The realised profit will be the payoff for the state that actually occurs.
\secInicio[sec:repair-shop-staffing][sec:summarized-solutions]

\clearpage
\sectionScored{Online Course Launch}{C7-\sml{3}}{sec:online-course-launch}
A training business must choose a full launch or a pilot launch before enrolment is known. Enrolment will be high ($S_1$), moderate ($S_2$), or low ($S_3$). Profits are in thousand USD.
\begin{center}\small\setlength{\tabcolsep}{4pt}
\begin{tabular}{@{}lrrr@{}}\toprule
Alternative&$S_1$&$S_2$&$S_3$\\\midrule
Full ($A_1$)&100&36&-8\\
Pilot ($A_2$)&52&48&28\\\bottomrule
\end{tabular}\end{center}
Suppose $P(S_1)=p$ and $P(S_2)=p/2$; determine $P(S_3)$ from the sum-to-one condition.
\QuestionsTasks[sec:online-course-launch]{
{7}/{sol:online-course-launch-c7}/{State the preferred alternative for every probability value ($0\le p\le1$). \score{3}}
}
\subsection*{C7 - Probability sensitivity analysis \score{3}}
\phantomsection\label{sol:online-course-launch-c7}
Alternatives: $A_1$: Full; $A_2$: Pilot.\par
States: $S_1$: High enrolment; $S_2$: Moderate enrolment; $S_3$: Low enrolment.
The states are mutually exclusive and exhaustive, so
\[
p+\frac p2+P(S_3)=1\quad\Longrightarrow\quad
P(S_3)=1-\frac{3p}{2},\qquad
\boldsymbol{q}=\left(p,\frac p2,1-\frac{3p}{2}\right).
\]
Every probability must lie in $[0,1]$:
\[
\begin{aligned}
0\le p\le1,&\qquad 0\le p/2\le1\ \Longrightarrow\ 0\le p\le2,\\
0\le1-3p/2\le1&\ \Longrightarrow\ 0\le p\le\frac23.
\end{aligned}
\]
Their intersection is $\boxed{0\le p\le\frac23}$.
Although the question asks about $[0,1]$, the interval
$\frac23<p\le1$ is infeasible because $P(S_3)<0$.
No decision based on these probabilities is valid there.
Using $E_i=EMV(A_i)=\sum_j P(S_j)P(A_i,S_j)$, in thousand USD:
\begin{align*}
E_1&=100p+36\left(\frac p2\right)-8\left(1-\frac{3p}{2}\right)
=100p+18p-8+12p=130p-8,\\
E_2&=52p+48\left(\frac p2\right)+28\left(1-\frac{3p}{2}\right)
=52p+24p+28-42p=28+34p.
\end{align*}
The threshold solves
\[
E_1=E_2\iff130p-8=28+34p\iff96p=36
\iff\boxed{p=\frac38}.
\]
It is feasible: $0<\frac38<\frac23$. At the tie, both EMVs equal
$\frac{163}{4}=40.75$ thousand USD.
The difference $E_1-E_2=96p-36$ is negative below $\frac38$ and positive above it.
At the feasible endpoints, $(E_1,E_2)=(-8,28)$ for $p=0$ and
$(E_1,E_2)=(\frac{236}{3},\frac{152}{3})$ for $p=\frac23$.
Thus both endpoints agree with the decision rule:
\[
\boxed{\begin{cases}
\text{Pilot launch }(A_2),&0\le p<\frac38,\\
\text{Full and pilot launches tied},&p=\frac38,\\
\text{Full launch }(A_1),&\frac38<p\le\frac23.
\end{cases}}
\]
Choose the pilot launch below the threshold and the full launch above it;
either launch is optimal at the tie. These are expected-profit decisions, not guarantees.
\secInicio[sec:online-course-launch][sec:summarized-solutions]

\clearpage
\sectionScored{Printing Contract}{C7-\sml{3}}{sec:printing-contract}
A printer must choose a standard or a flexible contract before a client's order schedule is known. The schedule will involve an early order ($S_1$), a revised order ($S_2$), or a late order ($S_3$). Profits are in thousand USD.
\begin{center}\small\setlength{\tabcolsep}{4pt}
\begin{tabular}{@{}lrrr@{}}\toprule
Alternative&$S_1$&$S_2$&$S_3$\\\midrule
Standard ($A_1$)&24&60&84\\
Flexible ($A_2$)&78&42&36\\\bottomrule
\end{tabular}\end{center}
Suppose $P(S_1)=2p$ and $P(S_2)=p$; determine $P(S_3)$ from the sum-to-one condition.
\QuestionsTasks[sec:printing-contract]{
{7}/{sol:printing-contract-c7}/{State the preferred alternative for every probability value ($0\le p\le1$).. \score{3}}
}
\subsection*{C7 - Probability sensitivity analysis \score{3}}
\phantomsection\label{sol:printing-contract-c7}
Alternatives: $A_1$: Standard; $A_2$: Flexible.\par
States: $S_1$: Early order; $S_2$: Revised order; $S_3$: Late order.
The three state probabilities sum to one:
\[
2p+p+P(S_3)=1\quad\Longrightarrow\quad
P(S_3)=1-3p,\qquad \boldsymbol{q}=(2p,p,1-3p).
\]
Requiring every probability to lie in $[0,1]$ gives
\[
0\le2p\le1\iff0\le p\le\frac12,\quad
0\le p\le1,\quad
0\le1-3p\le1\iff0\le p\le\frac13.
\]
The intersection is $\boxed{0\le p\le\frac13}$.
Although the question asks about $[0,1]$, the interval
$\frac13<p\le1$ is infeasible because $P(S_3)<0$.
For $p>\frac12$, $P(S_1)>1$ as well. No alternative is recommended in the infeasible interval.
The expected profits, in thousand USD, are
\begin{align*}
E_1&=EMV(A_1)=24(2p)+60p+84(1-3p)\\
&=48p+60p+84-252p=84-144p,\\
E_2&=EMV(A_2)=78(2p)+42p+36(1-3p)\\
&=156p+42p+36-108p=36+90p.
\end{align*}
Set the expected profits equal:
\[
E_1=E_2\iff84-144p=36+90p\iff48=234p
\iff\boxed{p=\frac8{39}}.
\]
The threshold lies within the domain: $0<\frac8{39}<\frac13$.
Both EMVs are $\frac{708}{13}$ thousand USD at the tie.
Since $E_1-E_2=48-234p$, the standard contract is better below $\frac8{39}$,
and the flexible contract is better above it.
At $p=0$, $(E_1,E_2)=(84,36)$; at $p=\frac13$, $(E_1,E_2)=(36,66)$.
These checks include both feasible endpoints.
\[
\boxed{\begin{cases}
\text{Standard contract }(A_1),&0\le p<\frac8{39},\\
\text{Standard and flexible contracts tied},&p=\frac8{39},\\
\text{Flexible contract }(A_2),&\frac8{39}<p\le\frac13.
\end{cases}}
\]
Choose the standard contract below the threshold and the flexible contract above it;
either contract is optimal at the tie. Expected profit is not guaranteed profit.
\secInicio[sec:printing-contract][sec:summarized-solutions]

\clearpage
\sectionScored{Summarized Solutions}{C1--C7}{sec:summarized-solutions}
{\small
The numbered summaries follow the exam order. $A_i$ and $S_j$ use the original
row and column order. $E_i=EMV(A_i)$; $b_j$ is the best payoff in state $S_j$;
$r_{ij}=b_j-P(A_i,S_j)$. Units are stated for each problem.
\subsection{Craft Workshop \hfill C1--3 C2--6}
\label{summary:craft-workshop}
Alternatives: $A_1$: Pottery workshop; $A_2$: Weaving workshop.\par
States: $S_1$: Strong demand; $S_2$: Weak demand.
Uncertain event: actual demand after the choice. Context: craft studio selecting a workshop; consequences are net contributions in USD.
\[P_{ij}=q_{ij}(v_i+d_j)+f_i.\]
The fixed sponsorship is added once; the two per-participant rates multiply attendance.
\begin{center}
\begin{tabular}{@{}lcc@{}}\toprule
Payoff (USD) & $S_1$: strong & $S_2$: weak\\\midrule
Pottery ($A_1$) & $30(12+6)+80=620$ & $18(12+2)+80=332$\\
Weaving ($A_2$) & $24(15+6)+50=554$ & $16(15+2)+50=322$\\
\bottomrule\end{tabular}
\end{center}
Pottery gives the higher contribution in both demand states.
\secInicio[sec:craft-workshop][sec:summarized-solutions]
\subsection{Guided Cycling Tours \hfill C1--3 C3--6}
\label{summary:guided-cycling-tours}
Alternatives: $A_1$: Town route; $A_2$: Countryside route.\par
States: $S_1$: Strong demand; $S_2$: Weak demand.
Uncertain event: actual demand after the choice. Context: tour operator selecting a route; consequences are seasonal profits in USD.
\[P_{ij}=q_{ij}(b_i+s_j)-q_{ij}c_i-F_j.\]
Revenue includes the surcharge; the seasonal fixed cost is subtracted once.
\begin{center}
\begin{tabular}{@{}lcc@{}}\toprule
Profit (USD) & $S_1$: strong & $S_2$: weak\\\midrule
Town ($A_1$) & $90(32+8)-90(18)-500=1480$ & $60(32+3)-60(18)-650=370$\\
Countryside ($A_2$) & $70(45+8)-70(26)-500=1390$ & $40(45+3)-40(26)-650=230$\\
\bottomrule\end{tabular}
\end{center}
The town route earns more in both demand states.
\secInicio[sec:guided-cycling-tours][sec:summarized-solutions]
\clearpage
\subsection{Sportswear Stock \hfill C4--2}
\label{summary:sportswear-stock}
Alternatives: $A_1$: Bulk order; $A_2$: Regular order; $A_3$: Small order; $A_4$: Pre-orders only.\par
States: $S_1$: High; $S_2$: Medium; $S_3$: Low.
All profits are in thousand USD.
\begin{center}
\begin{tabular}{@{}lcc@{}}\toprule
Alternative & Row maximum $M_i$ & Row minimum $m_i$\\\midrule
Bulk order ($A_1$) & 120 & -30\\
Regular order ($A_2$) & 84 & 16\\
Small order ($A_3$) & 55 & 29\\
Pre-orders only ($A_4$) & 24 & 24\\
\bottomrule\end{tabular}
\end{center}
\[\max_i M_i=\max\{120,84,55,24\}=120,\qquad \max_i m_i=\max\{-30,16,29,24\}=29.\]
\textbf{Maximax:} choose \textbf{Bulk order} ($A_1$), with best-case profit 120 thousand USD.\par
\textbf{Maximin:} choose \textbf{Small order} ($A_3$), with worst-case profit 29 thousand USD.
\secInicio[sec:sportswear-stock][sec:summarized-solutions]
\subsection{Concert Venue \hfill C4--2}
\label{summary:concert-venue}
Alternatives: $A_1$: Stadium; $A_2$: Theatre; $A_3$: Community hall; $A_4$: Online event.\par
States: $S_1$: High; $S_2$: Medium; $S_3$: Low.
All profits are in thousand USD.
\begin{center}
\begin{tabular}{@{}lcc@{}}\toprule
Alternative & Row maximum $M_i$ & Row minimum $m_i$\\\midrule
Stadium ($A_1$) & 150 & -45\\
Theatre ($A_2$) & 95 & 12\\
Community hall ($A_3$) & 60 & 26\\
Online event ($A_4$) & 30 & 30\\
\bottomrule\end{tabular}
\end{center}
\[\max_i M_i=\max\{150,95,60,30\}=150,\qquad \max_i m_i=\max\{-45,12,26,30\}=30.\]
\textbf{Maximax:} choose \textbf{Stadium} ($A_1$), with best-case profit 150 thousand USD.\par
\textbf{Maximin:} choose \textbf{Online event} ($A_4$), with worst-case profit 30 thousand USD.
\secInicio[sec:concert-venue][sec:summarized-solutions]
\subsection{Cold Storage \hfill C4--2}
\label{summary:cold-storage}
Alternatives: $A_1$: Build large; $A_2$: Build medium; $A_3$: Rent small; $A_4$: Shared storage.\par
States: $S_1$: Very strong sales; $S_2$: Strong sales; $S_3$: Steady sales; $S_4$: Weak sales.
All profits are in thousand USD.
\begin{center}
\begin{tabular}{@{}lcc@{}}\toprule
Alternative & Row maximum $M_i$ & Row minimum $m_i$\\\midrule
Build large ($A_1$) & 140 & -40\\
Build medium ($A_2$) & 100 & 8\\
Rent small ($A_3$) & 68 & 22\\
Shared storage ($A_4$) & 48 & 28\\
\bottomrule\end{tabular}
\end{center}
\[\max_i M_i=\max\{140,100,68,48\}=140,\qquad \max_i m_i=\max\{-40,8,22,28\}=28.\]
\textbf{Maximax:} choose \textbf{Build large} ($A_1$), with best-case profit 140 thousand USD.\par
\textbf{Maximin:} choose \textbf{Shared storage} ($A_4$), with worst-case profit 28 thousand USD.
\secInicio[sec:cold-storage][sec:summarized-solutions]
\clearpage
\subsection{Flower Stall \hfill C5--2}
\label{summary:flower-stall}
Alternatives: $A_1$: Large purchase; $A_2$: Medium purchase; $A_3$: Small purchase.\par
States: $S_1$: High; $S_2$: Medium; $S_3$: Low.
\[\boldsymbol{b}=\bigl(78,48,30\bigr),\qquad b_j=\max_k P(A_k,S_j).\]
All payoffs and regrets in this calculation are in thousand USD.
\begin{center}
\begin{tabular}{@{}lcccc@{}}\toprule
Alternative & $S_1$ & $S_2$ & $S_3$ & Maximum regret\\\midrule
$A_1$ & $\begin{gathered}78-78\\=0\end{gathered}$ & $\begin{gathered}48-32\\=16\end{gathered}$ & $\begin{gathered}30-(-12)\\=42\end{gathered}$ & 42\\
$A_2$ & $\begin{gathered}78-60\\=18\end{gathered}$ & $\begin{gathered}48-48\\=0\end{gathered}$ & $\begin{gathered}30-18\\=12\end{gathered}$ & 18\\
$A_3$ & $\begin{gathered}78-40\\=38\end{gathered}$ & $\begin{gathered}48-38\\=10\end{gathered}$ & $\begin{gathered}30-30\\=0\end{gathered}$ & 38\\
\bottomrule\end{tabular}
\end{center}
\[\min_i\max_j r(A_i,S_j)=\min\{42,18,38\}=18.\]
Choose \textbf{Medium purchase} ($A_2$): its maximum regret is \textbf{18 thousand USD}. 
\secInicio[sec:flower-stall][sec:summarized-solutions]
\subsection{Parcel Depot \hfill C5--2}
\label{summary:parcel-depot}
Alternatives: $A_1$: New depot; $A_2$: Expand depot; $A_3$: Rent space; $A_4$: Current depot.\par
States: $S_1$: High; $S_2$: Normal; $S_3$: Low.
\[\boldsymbol{b}=\bigl(112,64,36\bigr),\qquad b_j=\max_k P(A_k,S_j).\]
All payoffs and regrets in this calculation are in thousand USD.
\begin{center}
\begin{tabular}{@{}lcccc@{}}\toprule
Alternative & $S_1$ & $S_2$ & $S_3$ & Maximum regret\\\midrule
$A_1$ & $\begin{gathered}112-112\\=0\end{gathered}$ & $\begin{gathered}64-40\\=24\end{gathered}$ & $\begin{gathered}36-(-28)\\=64\end{gathered}$ & 64\\
$A_2$ & $\begin{gathered}112-86\\=26\end{gathered}$ & $\begin{gathered}64-64\\=0\end{gathered}$ & $\begin{gathered}36-14\\=22\end{gathered}$ & 26\\
$A_3$ & $\begin{gathered}112-62\\=50\end{gathered}$ & $\begin{gathered}64-54\\=10\end{gathered}$ & $\begin{gathered}36-34\\=2\end{gathered}$ & 50\\
$A_4$ & $\begin{gathered}112-36\\=76\end{gathered}$ & $\begin{gathered}64-36\\=28\end{gathered}$ & $\begin{gathered}36-36\\=0\end{gathered}$ & 76\\
\bottomrule\end{tabular}
\end{center}
\[\min_i\max_j r(A_i,S_j)=\min\{64,26,50,76\}=26.\]
Choose \textbf{Expand depot} ($A_2$): its maximum regret is \textbf{26 thousand USD}. 
\secInicio[sec:parcel-depot][sec:summarized-solutions]
\clearpage
\subsection{Bakery Output \hfill C5--2}
\label{summary:bakery-output}
Alternatives: $A_1$: High output; $A_2$: Flexible output; $A_3$: Low output.\par
States: $S_1$: Very high; $S_2$: High; $S_3$: Moderate; $S_4$: Low.
\[\boldsymbol{b}=\bigl(960,660,460,360\bigr),\qquad b_j=\max_k P(A_k,S_j).\]
All payoffs and regrets in this calculation are in USD.
\begin{center}
\begin{tabular}{@{}lccccc@{}}\toprule
Alternative & $S_1$ & $S_2$ & $S_3$ & $S_4$ & Maximum regret\\\midrule
$A_1$ & $\begin{gathered}960-960\\=0\end{gathered}$ & $\begin{gathered}660-560\\=100\end{gathered}$ & $\begin{gathered}460-180\\=280\end{gathered}$ & $\begin{gathered}360-(-240)\\=600\end{gathered}$ & 600\\
$A_2$ & $\begin{gathered}960-740\\=220\end{gathered}$ & $\begin{gathered}660-660\\=0\end{gathered}$ & $\begin{gathered}460-460\\=0\end{gathered}$ & $\begin{gathered}360-120\\=240\end{gathered}$ & 240\\
$A_3$ & $\begin{gathered}960-480\\=480\end{gathered}$ & $\begin{gathered}660-500\\=160\end{gathered}$ & $\begin{gathered}460-440\\=20\end{gathered}$ & $\begin{gathered}360-360\\=0\end{gathered}$ & 480\\
\bottomrule\end{tabular}
\end{center}
\[\min_i\max_j r(A_i,S_j)=\min\{600,240,480\}=240.\]
Choose \textbf{Flexible output} ($A_2$): its maximum regret is \textbf{240 USD}. 
\secInicio[sec:bakery-output][sec:summarized-solutions]
\clearpage
\subsection{Fruit Packing \hfill C6--2}
\label{summary:fruit-packing}
Alternatives: $A_1$: Premium packs; $A_2$: Mixed packs; $A_3$: Basic packs.\par
States: $S_1$: Strong; $S_2$: Steady; $S_3$: Weak.
All EMVs are in thousand USD.
\[\boldsymbol{q}=(0.40,0.35,0.25),\qquad 0.40+0.35+0.25=1.\]
All probabilities lie in $[0,1]$, so they are nonnegative and valid.
\begin{align*}
EMV(A_1)&=0.40(90)+0.35(44)+0.25(-12)=48.40,\\
EMV(A_2)&=0.40(68)+0.35(58)+0.25(24)=53.50,\\
EMV(A_3)&=0.40(46)+0.35(48)+0.25(38)=44.70.
\end{align*}
\[\max\{48.40,53.50,44.70\}=53.50.\]
Choose \textbf{Mixed packs} ($A_2$), with expected profit \textbf{53.50 thousand USD} (USD 53,500). 
Expected profit is not guaranteed profit.
\secInicio[sec:fruit-packing][sec:summarized-solutions]
\subsection{Art Market Space \hfill C6--2}
\label{summary:art-market-space}
Alternatives: $A_1$: Large tent; $A_2$: Medium tent; $A_3$: Small tent; $A_4$: Existing rooms.\par
States: $S_1$: High; $S_2$: Normal; $S_3$: Low.
All EMVs are in thousand USD.
\[\boldsymbol{q}=(0.30,0.40,0.30),\qquad 0.30+0.40+0.30=1.\]
All probabilities lie in $[0,1]$, so they are nonnegative and valid.
\begin{align*}
EMV(A_1)&=0.30(104)+0.40(38)+0.30(-24)=39.20,\\
EMV(A_2)&=0.30(78)+0.40(56)+0.30(18)=51.20,\\
EMV(A_3)&=0.30(56)+0.40(50)+0.30(36)=47.60,\\
EMV(A_4)&=0.30(28)+0.40(28)+0.30(28)=28.
\end{align*}
\[\max\{39.20,51.20,47.60,28\}=51.20.\]
Choose \textbf{Medium tent} ($A_2$), with expected profit \textbf{51.20 thousand USD} (USD 51,200). 
Expected profit is not guaranteed profit.
\secInicio[sec:art-market-space][sec:summarized-solutions]
\subsection{Repair Shop Staffing \hfill C6--2}
\label{summary:repair-shop-staffing}
Alternatives: $A_1$: Hire full team; $A_2$: Hire part-time; $A_3$: Current team.\par
States: $S_1$: Surge; $S_2$: Strong; $S_3$: Steady; $S_4$: Slow.
All EMVs are in thousand USD.
\[\boldsymbol{q}=(0.20,0.30,0.30,0.20),\qquad 0.20+0.30+0.30+0.20=1.\]
All probabilities lie in $[0,1]$, so they are nonnegative and valid.
\begin{align*}
EMV(A_1)&=0.20(120)+0.30(72)+0.30(24)+0.20(-36)=45.60,\\
EMV(A_2)&=0.20(88)+0.30(76)+0.30(52)+0.20(12)=58.40,\\
EMV(A_3)&=0.20(58)+0.30(60)+0.30(54)+0.20(44)=54.60.
\end{align*}
\[\max\{45.60,58.40,54.60\}=58.40.\]
Choose \textbf{Hire part-time} ($A_2$), with expected profit \textbf{58.40 thousand USD} (USD 58,400). 
Expected profit is not guaranteed profit.
\secInicio[sec:repair-shop-staffing][sec:summarized-solutions]
\clearpage
\subsection{Online Course Launch \hfill C7--3}
\label{summary:online-course-launch}
$A_1$: Full launch; $A_2$: Pilot launch. States: high, moderate, low enrolment.
Profits and EMVs are in thousand USD.
\[
P(S_3)=1-p-\tfrac p2=1-\tfrac{3p}{2},\qquad
\boldsymbol{q}=\left(p,\tfrac p2,1-\tfrac{3p}{2}\right).
\]
The conditions $0\le p\le1$, $0\le p/2\le1$ and $0\le1-3p/2\le1$
intersect in $0\le p\le\frac23$. The remaining interval $(\frac23,1]$ gives a negative third probability.
\begin{align*}
E_1&=100p+36(p/2)-8(1-3p/2)=130p-8,\\
E_2&=52p+48(p/2)+28(1-3p/2)=28+34p.\\
E_1=E_2&\iff130p-8=28+34p\iff96p=36\iff p=\tfrac38.
\end{align*}
The threshold is feasible. $E_1-E_2=96p-36$ changes from negative to positive.
At $p=0$, $(E_1,E_2)=(-8,28)$; at $p=\frac23$, $(E_1,E_2)=(\frac{236}{3},\frac{152}{3})$.
\[
\boxed{\begin{cases}
\text{Pilot launch }(A_2),&0\le p<\frac38,\\
\text{Full and pilot launches tied},&p=\frac38,\\
\text{Full launch }(A_1),&\frac38<p\le\frac23.
\end{cases}}
\]
No recommendation is valid for $\frac23<p\le1$.
\secInicio[sec:online-course-launch][sec:summarized-solutions]
\clearpage
\subsection{Printing Contract \hfill C7--3}
\label{summary:printing-contract}
$A_1$: Standard contract; $A_2$: Flexible contract. States: early, revised, late order.
Profits and EMVs are in thousand USD.
\[
P(S_3)=1-2p-p=1-3p,\qquad \boldsymbol{q}=(2p,p,1-3p).
\]
The restrictions $0\le2p\le1$, $0\le p\le1$ and $0\le1-3p\le1$
intersect in $0\le p\le\frac13$. The interval $(\frac13,1]$ gives $P(S_3)<0$;
also $P(S_1)>1$ when $p>\frac12$.
\begin{align*}
E_1&=24(2p)+60p+84(1-3p)=84-144p,\\
E_2&=78(2p)+42p+36(1-3p)=36+90p.\\
E_1=E_2&\iff84-144p=36+90p\iff48=234p\iff p=\tfrac8{39}.
\end{align*}
The threshold is feasible. $E_1-E_2=48-234p$ changes from positive to negative.
At $p=0$, $(E_1,E_2)=(84,36)$; at $p=\frac13$, $(E_1,E_2)=(36,66)$.
\[
\boxed{\begin{cases}
\text{Standard contract }(A_1),&0\le p<\frac8{39},\\
\text{Standard and flexible contracts tied},&p=\frac8{39},\\
\text{Flexible contract }(A_2),&\frac8{39}<p\le\frac13.
\end{cases}}
\]
No recommendation is valid for $\frac13<p\le1$.
\secInicio[sec:printing-contract][sec:summarized-solutions]
}
\end{document}
